\section{Two real quartic reflection-field class certificates at conductor sixteen}
\label{sec:conductor16-real-quartic-class}

Put $p=269$, $K=\mathbb Q(\sqrt{-3})$, $t=\sqrt[3]p>0$,
and $d=\sqrt{2p}>0$.  Define
\[
 \begin{gathered}
 \beta_1=\sqrt{2p+23d},\qquad \beta_2=\sqrt{2p+3d},\\
 F_i=\mathbb Q(\beta_i),\quad L_i=F_i(t),\quad M_i=KL_i.
 \end{gathered}
 \tag{G147C.1}
\]
Both inner radicands and their conjugates are positive.
Let $\phi$ be the primitive odd quartic character modulo sixteen with
$\phi(13)=i$, $\phi(5)=-i$, and let
$\psi=\omega_p^{67}$, where the Teichmuller character is normalized
by $i\equiv82\pmod p$.  Thus $\psi(2)=-i$.
Set
\[
 \xi_1=\phi\psi,\qquad \xi_2=\phi\psi^{-1},\qquad
 \delta=\xi_i^2=\chi_8\tau_p.
 \tag{G147C.2}
\]
The quartic characters are primitive even of conductor $16p=4304$;
$\delta$ is primitive even of fundamental conductor $8p=2152$.
Characters are zero at integers not coprime to their primitive conductor.

\begin{theorem}[AC-G147-class: the two real conductor-sixteen class sectors]
\label{AC-G147-class}
For both fields in (G147C.1),
\[
                      \operatorname{Cl}(L_i)[269]=0.
 \tag{G147C.3}
\]
Consequently the nontrivial cubic sector of
$\operatorname{Cl}(M_i)[269]$ vanishes.
For each $i$ the proof gives $h_{L_i}=14m_i$, where $m_i$ is the
index of seven explicitly certified units and $269\nmid m_i$.
It does not assert $m_i=1$, a complete class group, or the vanishing
of the trivial cubic class sector of $M_i$.
No strict-group or growing-window conclusion is included in this theorem.
\end{theorem}

\begin{proof}
\proofstep{1. Exact characters, fields, local data and signatures.}
Here is a finite cyclotomic construction fixing the character orientation.
The integer $2$ generates $\mathbb F_{269}^{\times}$.
In $\mathbb Q(z)$, $z=\zeta_p$, put
\[
 \eta_j=\sum_{a\in2^j(\mathbb F_p^\times)^4}z^a,\quad
 v=\eta_0-\eta_2,\quad v'=\eta_1-\eta_3,\quad
 s=\sum_{a=1}^{p-1}(a/p)z^a.
 \tag{G147C.4}
\]
Integer polynomial arithmetic modulo $\Phi_{269}$ gives
\[
 s^2=p,\qquad 2v^2=-p-13s,\qquad vv'=-5s,
 \qquad v^4+pv^2+25p=0.
 \tag{G147C.5}
\]
For reproducibility the exact verifier uses length-$269$ integer
coefficient vectors and cyclic convolution.  A vector vanishes at
$z$ exactly when all its coefficients are equal, since the only
rational relation of degree at most $268$ is $\Phi_{269}$.
It checks every identity in (G147C.5), primitivity of $2$, and
$z\mapsto z^2:v\mapsto v'\mapsto-v\mapsto-v'$.
In particular the displayed quartic lies in the conductor-$p$
cyclotomic field.  It has degree four: its square generates the
quadratic field and has nonsquare rational norm $25p$.
The corresponding character has value $\psi(2)=-i$ on this generator.
The finite cosine sum for $s$ is positive (checked with outward balls),
so $s=\sqrt p$ in the standard complex embedding.

Put $w_0=\zeta_{16}^{3}-\zeta_{16}^{-3}$ and $r_0=\sqrt2$.
The automorphism $b_0:\zeta_{16}\mapsto\zeta_{16}^{13}$ satisfies
\[
 w_0^2=-2-r_0,\quad b_0(w_0)=(r_0-1)w_0,\quad
 w_0b_0(w_0)=-r_0,\quad \phi(b_0)=i.
 \tag{G147C.6}
\]
The two cyclotomic quartic fields are disjoint, since their
conductors are coprime.  Set $r=w_0v$ and $r'=b_0(w_0)v'$.
Then $r^2+(r')^2=2p+13r_0s$ and $rr'=5r_0s$.
The sum $r+r'$ is fixed by the diagonal generator $(b_0,b_p)$,
where $b_p:z\mapsto z^2$; the difference $r-r'$ is fixed by
$(b_0,b_p^{-1})$.  These are exactly the kernels of $\xi_1$ and
$\xi_2$, respectively.  Their squares are $2p+23r_0s$ and
$2p+3r_0s$.  This proves the field identification in (G147C.1),
up to conjugation, and gives
\[
 \begin{gathered}
 P_1(X)=X^4-1076X^2+4842,\qquad
 P_2(X)=X^4-1076X^2+284602,\\
 D(F_i)=(16p)^2(8p)=2^{11}p^3.
 \end{gathered}
 \tag{G147C.7}
\]
Their quartic characters have exact conductor sixteen at two
because $\phi(9)=-1$, and exact conductor $p$ at $p$ because
$\psi$ is a nontrivial tame quartic character.  Their squares
are the primitive quadratic character of $\mathbb Q(\sqrt{2p})$.
The abelian conductor--discriminant formula proves the last identity.
The norms of the two squared generators are $18p$ and $1058p$,
both rational nonsquares, so each generator has degree four.
The two kernels are distinct; these fields are not identified.
For later orientation, the generator induced by $b_0$ has
$\xi_i(b_0)=i$ in both fields and
\[
 b_0(\beta_1)=\frac{\beta_1^3-1067\beta_1}{69},\qquad
 b_0(\beta_2)=\frac{\beta_2^3-547\beta_2}{69}.
 \tag{G147C.8}
\]
Indeed the products with $\beta_i$ are $3d$ and $23d$; applying
$b_0$ again negates $\beta_i$.

The cubic field $T=\mathbb Q(t)$ has discriminant $-3p^2$.
At $p$ its defining polynomial is Eisenstein and tame.
At three, the equation $z+3z^2+3z^3=-30$ has a Hensel solution;
then $-(1+3z)$ is a cube root of $269$ in $\mathbb Q_3$.
The remaining quadratic local factor is $\mathbb Q_3(\sqrt{-3})$,
of discriminant exponent one.  At two $T$ is unramified.
The polynomials $P_i$ are Eisenstein at $p$, so $F_{i,p}$ is
fully ramified of degree four.  Composing with the cubic local
field, whose ramification degree is coprime to four, gives
ramification index twelve and residue degree one.  It is tame;
the discriminant exponent is eleven.  At two and three the other
field is unramified, so finite-etale base change computes the
exponents $3\cdot11=33$ and $4\cdot1=4$.
No other prime ramifies.  Consequently
\[
 \begin{gathered}
 [L_i:\mathbb Q]=12,\quad(r_1,r_2)=(4,4),\\
 D:=D(L_i)=2^{33}3^4p^{11}
 =371321379996073401540039447603459391488.
 \end{gathered}
 \tag{G147C.9}
\]
The signature follows by combining four real embeddings of $F_i$
with the one real and two complex embeddings of $T$.
Each $L_i$ has a real embedding, so $|\mu_{L_i}|=2$.
The normal closure $K(t)$ has group $S_3$ and unique quadratic
subfield $K$, different from the unique quadratic subfield of $F_i$.
Thus $\operatorname{Gal}(M_i/\mathbb Q)=S_3\times C_4$,
and $L_i$ is fixed by a reflection in the $S_3$ factor.

\proofstep{2. Seven exact units and the original unit-lattice index.}
For each $i$ use the basis $t^a\beta_i^b$, $0\le a<3$, $0\le b<4$,
with coordinate index $4a+b$.  The corresponding exact file
\texttt{fieldi/all\_units\_exact\_certificate.json} contains seven
rational arrays and their monic integral degree-twelve polynomials,
all with constant coefficient $\pm1$.
The input hashes are
\texttt{f921b8e9...} for $i=1$ and \texttt{ecb3d30a...} for $i=2$;
the full hashes and complete arrays are in the attached evidence manifest.
The two exact multiplication matrices $T,W$ satisfy
\[
 T^3=pI,\quad W^4-4pW^2+C_iI=0,\quad TW=WT,
 \qquad (C_1,C_2)=(4842,284602).
 \tag{G147C.10}
\]
The raw generator is $x=\beta_i-t$.
The verifier checks its resultant and both inverse rational maps,
then evaluates every candidate unit polynomial at $W-T$.
The first column is precisely the supplied tensor array.
Its exact characteristic polynomial is monic integral with
constant term $\pm1$; an additional Horner evaluation is zero.
Thus the values and their inverses are algebraic integers,
so they are actual units independently of the candidate extraction.
No provisional maximal order or complete-unit claim is used.
The arrays have at most $71$ and $132$ digits in their numerators.

For the residue certificate use $\ell=206593$, $\ell-1=768p$,
primitive root $g=10$, and $h=g^{768}=49873\pmod\ell$.
Complete trial division and exact power tests verify these assertions.
The three simple roots of $t^3-p$ and four simple roots of $P_i$
give twelve embeddings $L_i\hookrightarrow\mathbb Q_\ell$ by Hensel.
All unit denominators are invertible modulo $\ell$.
For every such embedding, the residues of the seven units give
logarithms $c_j\in\mathbb F_p$ through $\bar u_j^{768}=h^{c_j}$.
The exact verified selected minors are
\[
 C^{(1)}=\begin{pmatrix}
263&187&233&85&59&233&165\\
6&36&187&85&165&59&81\\
6&233&82&85&233&81&59\\
263&82&36&85&81&165&233\\
263&187&233&214&263&85&50\\
6&36&187&214&50&263&140\\
6&233&82&214&85&140&263
\end{pmatrix},
 \tag{G147C.11}
\]
\[
 C^{(2)}=\begin{pmatrix}
6&85&177&100&129&107&177\\
263&85&169&177&140&177&22\\
263&85&100&92&140&232&107\\
6&85&92&169&129&22&232\\
6&214&177&100&135&180&56\\
263&214&169&177&134&56&224\\
263&214&100&92&134&78&180
\end{pmatrix}.
 \tag{G147C.12}
\]
Their determinants are $35056345116121834\equiv55\pmod p$
and $-1782298273092498\equiv181\pmod p$.
The complete twelve-row residue data are independently reproducible
from the rational arrays and modular powers.
The logarithms kill $p$th powers and the roots of unity.
Hence the seven units are independent modulo $p$ in the rank-seven
lattice $\mathcal O_{L_i}^\times/\mu_{L_i}$, have full rank, and their
finite index $m_i$ is prime to $p$.
Indeed a primitive integer relation would contradict the matrix
rank; the same mod-$p$ rank makes the lattice map invertible modulo $p$.

\proofstep{3. The real regulator and an exact denominator bound.}
Let $\rho=p^{1/3}$ and $\omega=e^{2\pi i/3}$, and put
$A_i=\sqrt{538+k_i\sqrt{538}}$, $B_i=\sqrt{538-k_i\sqrt{538}}$,
where $k_1=23$, $k_2=3$.
The four real places are $(t,\beta_i)=(\rho,w)$ for
$w=A_i,-A_i,B_i,-B_i$; the four complex-place representatives
are $(\rho\omega,w)$ with the same four choices of $w$.
Form the eight-by-seven matrix $n_v\log|\sigma_v(u_j)|$,
with $n_v=1$ at real places and $n_v=2$ at complex places.
Delete the last row and take the absolute determinant $R_{U,i}$.
The projected weighted-log convention gives exactly $R_{U,i}=m_iR_{L_i}$,
with no Euclidean hyperplane-covolume factor.
The exact arrays are evaluated with $5000$-bit outward balls;
each column sum contains zero, and the determinants satisfy
\[
 \begin{gathered}
 73197417311598<R_{U,1}<73197417311599,\\
 1584756866269325<R_{U,2}<1584756866269326.
 \end{gathered}
 \tag{G147C.13}
\]
No provisional regulator is used as an error bound.

The rational isolation needs no lower bound for the full regulator.
For $n=12$, $r_2=4$ define
\[
 U=\left\lceil\frac{n!4^{r_2}(\lfloor\sqrt D\rfloor+1)}
 {n^n3^{r_2}}\right\rceil,\qquad
 H=4^n\left\lceil\sqrt[3]{U^4}\right\rceil.
 \tag{G147C.14}
\]
Minkowski and $\pi>3$ give an integral representative of every class
of norm at most $U$.  Euler factors and an integral comparison give
\[
 h_{L_i}\le U^{4/3}\zeta_{L_i}(4/3)
 \le U^{4/3}\zeta(4/3)^{12}\le4^{12}U^{4/3}\le H,
 \tag{G147C.15}
\]
since $\zeta(4/3)\le1+\int_1^\infty x^{-4/3}dx=4$.
Here exact integer arithmetic gives
\[
 U=3271839082898095,\qquad H=8149090553217531846697418752.
 \tag{G147C.16}
\]
If $r=h_{L_i}/m_i>1$ and $|r-c|<1/H$ for an integer $c$,
then $m_i<h_{L_i}\le H$.  Unless $h_{L_i}=cm_i$, the nonzero
integer $h_{L_i}-cm_i$ would imply $|r-c|\ge1/m_i>1/H$.
This isolates the rational value without assuming $m_i=1$.

\proofstep{4. Artin factorization and the even finite Fourier formulas.}
Let $\varrho$ be the standard representation of the $S_3$ splitting
field of $X^3-p$.  The permutation representation for $L_i$ is
$(1\oplus\varrho)\otimes(1\oplus\delta\oplus\xi_i\oplus\bar\xi_i)$.
With all ramified inertia-invariant factors retained, Artin formalism gives
\[
 Z_i:=\operatorname*{res}_{s=1}\zeta_{L_i}(s)
 =L(1,\delta)|L(1,\xi_i)|^2
 L(1,\varrho)L(1,\varrho\delta)|L(1,\varrho\xi_i)|^2.
 \tag{G147C.17}
\]
The quadratic factor has its primitive conductor $8p$, whereas the
quartic factors have primitive conductor $16p$.
For a primitive even character $\theta$ of conductor $q$, put
\[
 S(\theta)=\sum_{a=1}^{q-1}\bar\theta(a)
                 \log(2\sin(\pi a/q)).
 \tag{G147C.18}
\]
For $0<r<1$ expand $-\log|1-re^{i\vartheta}|$
as the absolutely convergent cosine series.
Multiplication by $\bar\theta(a)$ and finite summation gives
$-\sum_a\bar\theta(a)\log|1-r\zeta_q^a|
=\mathfrak g(\bar\theta)\sum_{n\ge1}\theta(n)r^n/n$.
Take $r\uparrow1$ by the finite outer sum and Abel's theorem
for the convergent periodic-character Dirichlet series.
The primitive finite Fourier identity and Parseval give
$|\mathfrak g(\bar\theta)|^2=q$.  Therefore
\[
 L(1,\theta)=-S(\theta)/\mathfrak g(\bar\theta),\qquad
 |L(1,\xi_i)|^2=|S(\xi_i)|^2/(16p).
 \tag{G147C.19}
\]
No odd-character Bernoulli formula is used.
For the primitive real even character $\delta$, conjugate terms
make its Gauss sum real, and its square is $8p$.
The separate finite ball calculation of
$\sum_a\delta(a)\cos(2\pi a/(8p))$ is strictly positive
(and lies between $46$ and $47$).  Thus the sign is fixed and
\[
 L(1,\delta)=-S(\delta)/\sqrt{8p}>0.
 \tag{G147C.20}
\]
The last positivity is also a strict output of the ball calculation.
This uses no quadratic class-number or complete-unit assertion.

\proofstep{5. Exact rank-two coefficients and conductors.}
The coefficients $L(s,\varrho)=\sum a_n n^{-s}$ are multiplicative,
with $a_1=1$, $a_{3^v}=1$ and $a_{p^v}=0$ for $v\ge1$.
For primes $\ell\ne3,p$,
\[
 a_\ell=\#\{x\bmod\ell:x^3=p\}-1,\qquad
 a_{\ell^{v+1}}=a_\ell a_{\ell^v}-(-3/\ell)a_{\ell^{v-1}}.
 \tag{G147C.21}
\]
At $\ell\equiv2\pmod3$ these are one in even degree and zero
in odd degree.  At $\ell\equiv1\pmod3$, the exact cubic residue
power test decides between $a_\ell=2$ and $-1$.
The special factors follow from the local algebras in Step 1.
The twisted sequences are exactly $a_n\delta(n)$ and $a_n\xi_i(n)$.
At two their inertia invariants vanish; at $p$ cubic order three
combined with quadratic order two or quartic order four gives
orders six or twelve, hence again no inertia invariants.
Thus zero extension inserts the correct Euler factors.

Write $\varrho=\operatorname{Ind}_K^{\mathbb Q}\nu$.
The finite cubic Kummer character $\nu$ has conductor $(p)$:
it is tame and nontrivial at $p$, while $p$ is a cube at three.
At two, $K_2/\mathbb Q_2$ is unramified quadratic and
$N(1+4a)\equiv1+4\operatorname{Tr}(\bar a)\pmod8$,
$N(1+8a)\equiv1+8\operatorname{Tr}(\bar a)\pmod{16}$.
The residue trace $\mathbb F_4\to\mathbb F_2$ is onto.
Since $\chi_8(5)=-1$ and $\phi(9)=-1$, the norm pullbacks
have exact conductor exponents three and four.
The additional $p$-characters are unramified at two.
At $p$ the residue norm from the unramified quadratic $K_p$ is
onto; the norm pullbacks of the quadratic and quartic characters
have orders two and four.  Their products with $\nu$ have
orders six and twelve, still tame with conductor exponent one.
There is no cancellation, and all twists are unramified at three.
Consequently the primitive Hecke moduli and rational conductors are
\[
 f=p,8p,16p,\qquad Q=3f^2=217083,13893312,55573248.
 \tag{G147C.22}
\]
\proofstep{6. Primitive finite Fourier transform and theta phase.}
The imaginary quadratic Minkowski bound for $K$ is
$2\sqrt3/\pi<2$, so $h(K)=1$.
Each of the three primitive Hecke characters is represented by
a multiplicative character $\nu_f$ on $\mathcal O_K/f\mathcal O_K$,
zero at nonunits and trivial on the six global units.
For the cubic character the latter also follows from
$9\mid p+1$; for the Dirichlet pullbacks it follows from the norm
of a unit being one.  Put
\[
 a=\frac{2\pi}{\sqrt3f},\qquad
 \Theta(t)=\frac16\sum_{\alpha\in\mathcal O_K}
    \nu_f(\alpha)e^{-a|\alpha|^2t}
       =\sum_{n\ge1}b_ne^{-ant}.
 \tag{G147C.23}
\]
The six generators of every ideal give the second identity.

The additive pairing
\[
 (\alpha,\beta)\longmapsto
 \exp\left(2\pi i\operatorname{Tr}_{K/\mathbb Q}
                   \frac{\alpha\bar\beta}{f\sqrt{-3}}\right)
 \tag{G147C.24}
\]
is perfect, by the inverse different.  At a unit frequency the
finite Fourier sum is $G\bar\nu_f(\bar\beta)$ by substitution.
At a nonunit frequency at a factor $\mathfrak q^a$ of the modulus,
choose $u\in1+\mathfrak q^{a-1}$ on which the character is
nontrivial, using primitivity; when $a=1$, use a residue unit.
Choose it to be one at the other factors.  Multiplication by $u$
fixes the additive character but multiplies the sum by
$\nu_f(u)\ne1$, so that sum is zero.
Parseval on the $f^2$ residue classes gives $|G|^2=f^2$.
This includes the prime powers eight and sixteen.

The lattice $f\mathcal O_K$ has area $f^2\sqrt3/2$ and
Euclidean dual $2/(f\sqrt{-3})\mathcal O_K$.
For $b=2/\sqrt{-3}$,
$\operatorname{Re}(\alpha\overline{b\beta})
 =\operatorname{Tr}_{K/\mathbb Q}(-\alpha\bar\beta/\sqrt{-3})$.
Thus the two in this dual cancels the one-half in the real-part
formula; no invertibility of two modulo $f$ is asserted.
Poisson summation with the Fourier transform
$\pi(at)^{-1}\exp(-\pi^2|\xi|^2/(at))$ of the Gaussian gives
\[
 \Theta(t)=\varepsilon t^{-1}\overline{\Theta(1/t)},
          \qquad |\varepsilon|=1.
 \tag{G147C.25}
\]
Indeed the scalar is $1/(ft)$, the dual exponent is
$a|\beta|^2/t$, and conjugating $\beta$ changes
$\bar\nu_f(\bar\beta)$ to $\bar\nu_f(\beta)$ without changing
its norm.  The scalar Gauss sign is absorbed into $\varepsilon$.
Both constant terms vanish.  Mellin integration, first for
$\operatorname{Re}s>1$ and then by the exponential bounds at both
ends, identifies the integral with $\Gamma(s)a^{-s}L(s)$ at $s=1$.
Splitting at one and substituting $t\mapsto1/t$ gives the exact identity
\[
 L(1)=\sum_{n\ge1}\frac{b_ne^{-an}}n
       +\varepsilon a\sum_{n\ge1}\bar b_n E_1(an),
 \qquad E_1(x)=\int_1^\infty e^{-xt}\frac{dt}{t}.
 \tag{G147C.26}
\]
The theta ball at one excludes zero in all three computations.
Therefore its actual phase is
$\varepsilon=\Theta(1)/\overline{\Theta(1)}$, evaluated with
inclusion arithmetic.  It is one for the two real sequences.

\proofstep{7. Tails, exact rational isolation and reflection descent.}
From (G147C.21), $|b_n|\le d(n)\le n$.
Writing $z_a=e^{-a}$ and $H_N=e^{-a(N+1)}$, the omitted tails
in the first sum of (G147C.26), the second sum before its factor
$a$, and the theta sum have respective bounds
\[
 \frac{H_N}{1-z_a},\qquad \frac{H_N}{a(1-z_a)},\qquad
 \frac{H_N((N+1)-Nz_a)}{(1-z_a)^2}.
 \tag{G147C.27}
\]
The middle estimate uses $E_1(x)\le e^{-x}/x$.
The script adds the full tail radius to each real and imaginary
component, uses $256$-bit balls with cutoffs $12000,80000,160000$,
and retains the $5000$-bit regulator ball.
The finite logarithm sums in (G147C.18)--(G147C.20) are evaluated
at every unit residue; each sine ball is strictly positive.
The complete output records all conductors, precisions, weighted-log
rows, theta/phase balls, tails and final intervals.

The analytic class-number formula, with four real places,
four complex places and two roots of unity, gives
\[
 h_{L_i}R_{L_i}=\frac{2\sqrt D}{2^4(2\pi)^4}Z_i.
 \tag{G147C.28}
\]
The resulting strict comparisons, valid separately for both fields, are
\[
 r_i=\frac{h_{L_i}R_{L_i}}{R_{U,i}}=\frac{h_{L_i}}{m_i}>1,
 \qquad |r_i-14|<2\cdot10^{-51},\qquad 1/H>1.22\cdot10^{-28}.
 \tag{G147C.29}
\]
The actual inclusion comparisons use the full balls, not these rounded
bounds.  The selector $14$ enters only this last comparison; it is
not used in a regulator, field, coefficient or error estimate.
Step 3 proves $h_{L_i}=14m_i$ exactly.
Since $269\nmid14m_i$, (G147C.3) follows.

Let $s$ be the reflection fixing $L_i$ in $M_i$.
Norm and extension on ideal classes have composites $2$ and $1+s$.
Since the $269$-primary class group of $L_i$ is zero, $1+s$ vanishes
on the $269$-primary class group of $M_i$.
All conjugate reflections act by $-1$ there, so their products,
the cubic rotations, act by one.
Exactly the nontrivial cubic class sector therefore vanishes;
no claim about the other sector is made.
\end{proof}

\paragraph{Actual-label boundary.}
For the original finite labels $\chi^{-j}\phi^{-\sigma}$ the residual
Kummer quotient at $e\equiv66\pmod{268}$ has quartic factor
$\phi^{-\sigma}\psi^{-1}$; at $e\equiv200$ it has
$\phi^{-\sigma}\psi$.
The fields are $F_1,F_2$ in the order $(\sigma=+1,-1)$ for the
first progression and in the reverse order for the second.
These are residual splitting fields only: no $\psi$ character is
inserted into the old finite host-character list.
The exact local field $M_{i,p}$ has ramification index twelve,
residue degree two and one prime above $p$.
Any subsequent strict argument must use this local situation and
retain all rational-basis denominators; it is not supplied by the
class certificate or by the earlier two-place linear test.

\paragraph{Sources and computation scope.}
The retained analytic class-number formula and projected weighted-log
normalization are Sutherland, MIT 18.785 Lecture 19 (2019), pp.1 and 4,
\url{https://math.mit.edu/classes/18.785/2019fa/LectureNotes19.pdf}.
The classical Minkowski ideal bound is retained as in AC-K19;
the denominator isolation is proved in Step 3.
The Dirichlet even logarithm formula is proved here from finite
Fourier summation and Abel convergence, with an independently
verified sign for the real quadratic Gauss sum.
The conductor--discriminant formula, Artin formalism, Dirichlet's
unit theorem and ordinary Poisson summation retain their previously
audited scopes; no fresh complete proof of those deep inputs is claimed.
Every finite Gaussian-period, unit and residue identity is tested
by exact integer, rational or modular arithmetic in saved scripts.
The numerical runtime is python-flint 0.9.0 and documented outward
ball operations, including order-one \texttt{arb.expint}; its
implementation correctness is not independently reproved.
Candidate-only PARI extraction supplies rational unit candidates
and the integer selector.  Neither a provisional class group,
a complete unit basis, GRH nor \texttt{bnfcertify} is used.
Full hashes bind the exact evidence and executable inputs.
